% journal_intro.tex — Section 1 of the journal manuscript (paper.tex).

\section{Introduction}
Evaluating measurement uncertainty by the propagation of distributions,
following Supplement~1 to the Guide to the Expression of Uncertainty in
Measurement (GUM~S1)~\cite{jcgm100,jcgm101}, requires the measurement model
$Y=f(\mathbf X)$ to be evaluated once per Monte Carlo trial, typically
$10^5$ to $10^6$ times~\cite{cox2006,harris2014}. For a growing class of
quantities $f$ is a numerical simulation. A computational fluid dynamics
(CFD) model of an accidental ammonia release in a factory~\cite{tan2017}
can take hours to days per run, so that a practical budget is a few tens of
runs, three or four orders of magnitude short of what the Monte Carlo
method needs. The same gap arises in scatterometry, where it has been
studied in detail~\cite{heidenreich2017}, and in many inverse and
model-based measurements.

The usual remedy is a \emph{surrogate} (emulator, metamodel) fitted to the
available runs and evaluated in place of $f$~\cite{ohagan2006,wang2022}.
Two questions then decide whether the resulting uncertainty statement can
be trusted: which surrogate reproduces the reported quantities best for a
given number of runs, and how the surrogate's own error is to be included
in the uncertainty budget, since it is not negligible at such budgets.
This paper answers both with a controlled comparison on ten models,
including a CFD dispersion case, and validates the answers against the
true models.

\subsection{Related work}
\emph{Gaussian process emulators.} Gaussian process (GP) regression, or
kriging, is the standard surrogate for deterministic computer
experiments~\cite{sacks1989,rw2006}. The Bayesian treatment of O'Hagan and
co-workers~\cite{ohagan2006,kennedy2001} reads the GP posterior as the
uncertainty about $f$ itself, which suggests using the posterior variance
as the surrogate contribution to an uncertainty budget. A common guideline
asks for about ten runs per input~\cite{loeppky2009}, and the effect of
modelling choices such as the correlation family, the mean function and
the treatment of the hyperparameters on GP accuracy has been assessed
systematically~\cite{chen2016}.

\emph{Polynomial chaos.} Polynomial chaos expansions (PCE) represent $Y$ in
polynomials orthonormal with respect to $g_{\mathbf X}$~\cite{xiu2002},
from which the moments and sensitivity indices follow
directly~\cite{sudret2008}. With few runs, sparse expansions selected by
least-angle regression~\cite{efron2004} and a corrected leave-one-out
error are the method of choice~\cite{blatman2011}.

\emph{Ensembles and committees.} Combining several surrogates with weights
derived from their cross-validation errors protects against choosing a poor
surrogate~\cite{goel2007,viana2009}; the spread of the members has also
been used as an interval for the output, in the spirit of probability
boxes~\cite{ferson1996}.

\emph{Adaptive designs.} Instead of fixing all runs in advance, runs can be
added where the surrogate is most uncertain, by its posterior
variance~\cite{gramacy2009}, by the disagreement of a committee, or by
measures such as the spectral uncertainty of Lin et al.~\cite{lin2024}. For
the probability of exceeding a threshold, the AK-MCS method of Echard et
al.~\cite{echard2011} places runs where the sign of $f-T$ is most
uncertain.

\emph{Validity of the surrogate contribution.} Whether any of these
uncertainty descriptions has the coverage it claims is rarely checked
against the true model. Conformal prediction~\cite{angelopoulos2023}, and
in particular the jackknife+~\cite{barber2021}, gives prediction intervals
with a coverage guarantee that holds for any fitted model, using only
leave-one-out fits. To our knowledge it has not been compared with GP
posterior intervals as the surrogate contribution in a GUM~S1 evaluation.

\emph{Benchmarks.} Comparisons of surrogates are sensitive to the test
functions and to the number of repetitions. We use standard functions from
the literature~\cite{vlse,wicaksono2023}, GUM-type input distributions, a
physical dispersion formula and a CFD case, with 30 paired repetitions and
statistics adjusted for multiple comparisons.

\subsection{Earlier analysis and contributions}
An earlier analysis by the first author combined a heterogeneous committee,
the spectral acquisition criterion and dynamic model management, and
reported that a passively sampled GP fails catastrophically on large input
domains, a failure the committee avoided. The present study re-examines
these claims with correctly specified baselines and extends them to the
general question above. Its contributions are:
\begin{enumerate}\setlength\itemsep{1pt}
\item a paired comparison, on ten models with uniform and GUM-type inputs
and a CFD dispersion case, of GP specifications, sparse PCE, committees and
adaptive and targeted designs at budgets of 10 to 80 runs, scored on the
quantities a GUM~S1 evaluation reports;
\item the finding that the GP failure of the earlier analysis is a
kernel-specification artefact, and that specification (amplitude,
anisotropy, trend) and variable transformation matter more than the choice
between single surrogates, committees and adaptive designs;
\item a validation, against the true models, of three ways to include the
surrogate error in GUM~S1: the GP posterior, a cross-validation (CV)
normal error and resampled leave-one-out residuals, with jackknife+
intervals as a distribution-free reference;
\item a procedure for surrogate-based GUM~S1 evaluations
(Fig.~\ref{fig:flowchart}), and code, per-seed records and an audit script
that recompute every number reported~\cite{repo}.
\end{enumerate}
